\documentclass[10pt]{article}
\usepackage[margin=0.75in]{geometry}
\usepackage{booktabs}
\usepackage{amsmath}
\usepackage{hyperref}
\title{Causal-Temporal Machine Learning for Root-Cause Localization in 5G Networks:\\ A Benchmark Audit and a Cautionary Evaluation}
\author{Nagi Reddy SaiBaba (23BDS0019) \and Arya Diwakar (23BDS0023) \and Jeswanth G (23BDS0212)}
\date{\today}
\begin{document}
\maketitle

\begin{abstract}
We audit TelecomTS and find that our initial headline does not survive
correction of the benchmark. We repair a train/test background leak and
injection-label errors, expose a per-fault-type template shortcut, add
instance-varying injections that test it, and retrain on the repaired split (5,689 test
windows, 424 positives, 26 recordings). The prior result collapses. TAGN, our
temporal-attention graph network, reaches 0.7589 PR-AUC across five seeds,
below TCN (0.8606) and logistic regression (0.8196). Removing the graph raises
PR-AUC from 0.7439 to 0.7581: the graph is not a marginal contributor, it is
worse than no graph in a single-seed ablation. The learned graph estimates
statistical dependencies; it is not a verified causal graph. The intervention-based RCA metric fails its own positive control:
the ground-truth root set scores 0.496 $\pm$ 0.052, essentially the
random-root null at 0.503 $\pm$ 0.016; the full true set against a matched
random set gives 0.508 $\pm$ 0.085. The predicted-root score
(0.567 $\pm$ 0.235) is unstable across seeds, ranging from 0.304 to 0.976, not
evidence of localization. The old argmax-lag early-warning number is also an
artifact: under a calibrated event-level protocol, detection at a 1\%
false-alarm budget is 3.7\%, with per-step AUC 0.495. We claim no
state-of-the-art model. We claim a corrected benchmark, a validation protocol
that separates real RCA from benchmark artifacts, and a decisive negative
result: the graph provides no clear detection benefit in the ablation, the
causal interpretation fails validation, and our initial headline is unsupported.
\end{abstract}

\section{Audit and Repairs}
We audited TelecomTS along four axes: split integrity, preprocessing leakage,
ground-truth validity, and metric protocol. Three findings changed the numbers:
(G5) a train/test background leak in which a test anomaly's normal source was
chunk-split into train; (G2--G4) injection label errors (identity/inert effects
labeled, multi-event windows under-labeled, \texttt{end\_time} 10$\times$ too
large); and a template shortcut in which the affected-KPI set is constant per
fault type. An input-free type-oracle reaches 0.971 top-$k$ Jaccard on the
shipped data, above every trained model; after instance-varying injection the
ceiling drops to 0.661. We bind sources atomically to their injected twins,
regenerate the corpus (280 spurious labels removed across 237/714 windows),
and refit scalers and graph prototypes per leave-recordings-out fold. The
repaired test set has $n=5{,}689$, 424 positives, and 26 recordings. The
repairs are not cosmetic. They reverse the headline.

\section{Detection}
\begin{table}[h]
\centering
\begin{tabular}{lcc}
\toprule
Model & PR-AUC & ROC-AUC \\
\midrule
TCN & \textbf{0.8606} & 0.9438 \\
LogReg & 0.8196 & 0.9309 \\
TAGN (graph + summary, 5 seeds) & 0.7589 & \textbf{0.9543} \\
GRU & 0.7635 & 0.9399 \\
\bottomrule
\end{tabular}
\caption{Setting I, repaired split, recording-level cluster-bootstrap protocol.}
\end{table}

TAGN does not win. It trails TCN and logistic regression on mean PR-AUC. Its
best seed (1337, 0.8278) edges logistic regression but remains below TCN. Paired
cluster-bootstrap tests against LogReg are negative on 4/5 seeds and every 95\%
interval includes zero: no advantage is detectable, and TAGN loses on the point
estimate. The single-seed component ablation is also negative. Removing the graph gives
0.7581 PR versus 0.7439 for the full model, so the graph is not neutral;
no\_graph is the best variant. Message passing contributes nothing beyond
summary statistics.

\section{Interventional Fidelity Fails Its Positive Control}
We define a model-internal do-style perturbation by zeroing selected normalized
KPI channels and re-scoring the anomaly head. The fidelity is the fraction of
windows in which intervening on a candidate root set reduces the score more
than intervening on a matched random non-root set. A valid metric must pass a
positive control on the ground-truth affected-KPI set. It does not: across five
checkpoints and 424 anomalous test windows from 13 recordings, true roots score
0.496 $\pm$ 0.052 at $k=3$, 0.476 $\pm$ 0.086 at $k=5$, and
0.519 $\pm$ 0.137 at $k=7$, while random roots score
0.503 $\pm$ 0.016, 0.506 $\pm$ 0.026, and 0.495 $\pm$ 0.025. The full true set
against a matched random set gives 0.508 $\pm$ 0.085, so truncation is not the
reason for the failure. This is not a borderline result: the ground-truth-root
score is at chance and overlaps the matched random-root null. The predicted-root score is
0.567 $\pm$ 0.235, but it is unstable across seeds (0.304 to 0.976); one
checkpoint drives the mean. Because the positive control fails, we do not
interpret this as evidence of causal localization. The earlier pilot value of
0.465 is invalid: it used a single checkpoint without matched controls and is
not reproducible under the corrected implementation. The learned graph is a
statistical dependency graph, not a verified causal DAG; the intervention is a
model-internal perturbation, not a physical do-operation.

\section{Early Warning}
The usual argmax-lag score reports 65.2\% ``warning before onset,'' but the
argmax can occur after onset and is not calibrated to a false-alarm budget. We
instead choose a threshold on validation normal windows and require the first
pre-onset alarm, then report event-level detection and false alarms. Over five
seeds and 93 test events, a 1\% false-alarm budget yields 3.7\% detection
(0.018 false alarms per normal window); at 5\% it yields 10.1\% detection
(0.096 false alarms per window). Per-step AUC is 0.495 and AP is 0.006. The old
number is not an early-warning result. Under a calibrated protocol the hazard
head is near chance and is not usable as an early-warning system on this data.

\section{Conclusion}
The audit changes the conclusions. After repair, TAGN trails TCN and logistic
regression on mean PR-AUC, the single-seed graph ablation favors no graph, the intervention metric fails its positive control, and the
calibrated hazard head is near chance. The learned graph is a dependency
estimate, not a causal mechanism. We put intervention-based RCA on trial and it
failed; we replaced the broken early-warning metric and the head failed there
too. These are decisive audit findings. We release the pipeline, the repaired
split, the validation scripts, and the test suite so that others can reproduce
the audit and apply the same controls to their own claims. The initial headline
does not hold under the repaired benchmark and protocol.

\end{document}
